Para cada \(\varepsilon>0\), sea \(y_\varepsilon\) la solución de
\[
\varepsilon y_\varepsilon''+y_\varepsilon'=0,
\qquad y_\varepsilon(0)=a,\qquad y_\varepsilon'(0)=b,
\qquad t\ge0.
\]
Todas las variables son adimensionales; \(a\) y \(b\) son constantes reales que no dependen de \(\varepsilon\).

Obtenga \(y_\varepsilon\) y compárela con la solución de la ecuación que resulta al poner \(\varepsilon=0\), conservando el dato \(y(0)=a\). ¿Qué sucede con el otro dato inicial? Dé una cota de \(|y_\varepsilon(t)-a|\) válida simultáneamente para todo \(t\ge0\). Compare también las derivadas en \(t=0\), en un tiempo \(t>0\) fijo y en \(t=\varepsilon\). Decida si esas derivadas pueden aproximarse a la derivada de la solución reducida con un error que tienda a cero para todos los tiempos \(t\ge0\) a la vez, e indique si hay un caso excepcional. Base la respuesta en las fórmulas explícitas.
